\documentclass[12pt,a4paper]{article}

% Paquetes útiles
\usepackage[utf8]{inputenc}
\usepackage[spanish]{babel}
\usepackage{geometry}
\usepackage{longtable}
\usepackage{hyperref}
\usepackage{graphicx}
\usepackage{booktabs}
\usepackage{caption}
\usepackage{array}

\geometry{margin=2.5cm}

\title{System Requirements Document (SRD) \\ Sistema de Gestión Documental (SGD)}
\author{Nombre de la Organización}
\date{\today}

\begin{document}

\maketitle
\tableofcontents
\newpage

% ======================================================
\section{Introducción}

\subsection{Propósito del documento}
Este documento define los requisitos funcionales, técnicos y de gestión para el desarrollo del Sistema de Gestión Documental (SGD) conforme a las normas ISO 15489, 30301, 23081, 16175 y 26102. Su objetivo es garantizar la autenticidad, integridad, trazabilidad y disponibilidad de los documentos durante todo su ciclo de vida.

\subsection{Alcance}
El sistema permitirá:
\begin{itemize}
    \item Capturar, registrar y clasificar documentos electrónicos y digitalizados.
    \item Gestionar metadatos según ISO 23081.
    \item Controlar acceso, roles y trazabilidad de uso mediante microservicio de usuarios.
    \item Garantizar retención y disposición final conforme a TRD.
    \item Integrarse con otros sistemas mediante eventos RabbitMQ y API REST/gRPC.
\end{itemize}

\subsection{Definiciones, acrónimos y abreviaturas}
\begin{itemize}
    \item \textbf{SGD}: Sistema de Gestión Documental.
    \item \textbf{TRD}: Tabla de Retención Documental.
    \item \textbf{ISO}: Organización Internacional de Normalización.
    \item \textbf{API}: Interfaz de Programación de Aplicaciones.
\end{itemize}

% ======================================================
\section{Marco Normativo}
\begin{longtable}{p{3cm}p{7cm}p{5cm}}
\toprule
\textbf{Norma ISO} & \textbf{Título} & \textbf{Aplicación al sistema} \\
\midrule
ISO 15489-1:2016 & Información y documentación — Gestión de documentos & Define procesos y principios de gestión documental \\
ISO 30301:2019 & Sistema de gestión para documentos — Requisitos & Establece requisitos organizacionales y auditoría \\
ISO 23081 & Metadatos para gestión de documentos & Define los metadatos esenciales y su gestión \\
ISO 16175 & Requisitos funcionales para sistemas electrónicos & Determina funciones que debe cumplir el software \\
ISO 26102:2019 & Requisitos técnicos y de conformidad para software SGD & Define requisitos técnicos y de calidad del software \\
\bottomrule
\end{longtable}

% ======================================================
\section{Requisitos Funcionales}

\subsection{Captura y Registro de Documentos}
\begin{itemize}
    \item Registro manual o automático desde correo, escáner, APIs.
    \item Asignación de identificador único (UUID).
    \item Captura de metadatos mínimos obligatorios.
\end{itemize}

\subsection{Clasificación Documental}
\begin{itemize}
    \item Aplicación de serie y subserie documental (TRD).
    \item Clasificaciones múltiples: proceso, tipo, proyecto.
    \item Registro de relaciones entre documentos (versiones, anexos).
\end{itemize}

\subsection{Control de Versiones}
\begin{itemize}
    \item Historial completo de versiones.
    \item Registro de autor, fecha y comentario por modificación.
    \item Restauración de versiones anteriores sin pérdida de integridad.
\end{itemize}

\subsection{Acceso y Seguridad}
\begin{itemize}
    \item Control de acceso basado en microservicio de usuarios.
    \item Auditoría de acciones con hash encadenado (integridad ISO 16175).
    \item Gestión de permisos y confidencialidad.
\end{itemize}

\subsection{Búsqueda e Indexación}
\begin{itemize}
    \item Búsqueda por texto libre y metadatos.
    \item Filtros por fecha, tipo documental, palabras clave.
    \item Integración con motor OpenSearch para índices optimizados.
\end{itemize}

\subsection{Retención y Disposición Final}
\begin{itemize}
    \item Integración con TRD configurable.
    \item Notificación de documentos próximos a disposición final.
    \item Registro de eliminación o transferencia con trazabilidad.
\end{itemize}

\subsection{Auditoría y Trazabilidad}
\begin{itemize}
    \item Registro de todas las acciones sobre documentos.
    \item Exportación de auditorías en PDF/A o XML.
    \item Encadenamiento de hash para integridad.
\end{itemize}

\subsection{Interoperabilidad}
\begin{itemize}
    \item Importación/exportación de documentos y metadatos (JSON, CSV, XML, PDF/A).
    \item Integración con APIs REST/gRPC.
    \item Cumplimiento de firma electrónica XAdES/PAdES.
\end{itemize}

% ======================================================
\section{Requisitos Técnicos}

\subsection{Arquitectura General}
\begin{itemize}
    \item Arquitectura modular Rust + Python.
    \item Soporte para on-premise y despliegue en contenedores.
    \item Base de datos PostgreSQL con metadatos JSONB.
\end{itemize}

\subsection{Metadatos}
\begin{longtable}{p{3cm}p{5cm}p{4cm}p{3cm}}
\toprule
\textbf{Categoría} & \textbf{Campo} & \textbf{Descripción} & \textbf{Obligatorio} \\
\midrule
Identificación & ID único & Identificador inalterable del documento & Sí \\
Contexto & Productor / Unidad & Área responsable & Sí \\
Contenido & Asunto / Título & Nombre del documento & Sí \\
Fecha & Fecha de creación & Timestamp del registro & Sí \\
Estructura & Formato & Tipo de archivo (PDF, DOCX, etc.) & Sí \\
Relación & Documento relacionado & ID o referencia cruzada & No \\
Gestión & Estado & Borrador, vigente, archivado & Sí \\
Seguridad & Nivel de acceso & Público, interno, confidencial & Sí \\
\bottomrule
\end{longtable}

\subsection{Seguridad y Rendimiento}
\begin{itemize}
    \item Cifrado TLS 1.3 y AES-256 en reposo.
    \item Autenticación multifactor.
    \item Escalabilidad mínima: 10.000 usuarios concurrentes.
    \item Búsqueda <2 segundos por consulta.
\end{itemize}

\subsection{Conformidad y Validación}
\begin{itemize}
    \item Documentación de pruebas unitarias y funcionales.
    \item Trazabilidad entre requisitos ISO y funcionalidades.
\end{itemize}

% ======================================================
\section{Requisitos de Gestión Documental}
\begin{itemize}
    \item Registro de políticas de gestión documental.
    \item Roles: Administrador, Productor, Usuario final, Auditor.
    \item Flujo de aprobación y validación documental.
    \item Indicadores: documentos creados, tiempos de acceso, cumplimiento de TRD.
\end{itemize}

% ======================================================
\section{Esquema de Metadatos y Crosswalk}
\begin{longtable}{p{3cm}p{3cm}p{3cm}p{3cm}p{2cm}}
\toprule
\textbf{Campo ISO 23081} & \textbf{Dublin Core} & \textbf{Schema.org} & \textbf{Tipo} & \textbf{Req} \\
\midrule
Identificador & Identifier & identifier & Texto & Sí \\
Título / Asunto & Title & name & Texto & Sí \\
Productor & Creator & creator & Texto & Sí \\
Fecha & Date & dateCreated & Fecha & Sí \\
Tipo documental & Type & genre & Texto & Sí \\
Palabras clave & Subject & keywords & Texto & No \\
Descripción & Description & description & Texto & No \\
Estado & Status & creativeWorkStatus & Enumerado & Sí \\
Nivel de acceso & AccessRights & accessMode & Enumerado & Sí \\
\bottomrule
\end{longtable}

% ======================================================
\section{Anexos}

\subsection{Glosario}
\begin{itemize}
    \item Documento: Información registrada como evidencia de actividades.
    \item Metadato: Dato que describe contexto, contenido y estructura de un documento.
    \item Serie documental: Conjunto de documentos relacionados por función o actividad.
    \item Autenticidad: Garantía de que el documento es lo que pretende ser.
    \item Integridad: Que el documento no ha sido alterado de manera no autorizada.
\end{itemize}

\subsection{Matriz de trazabilidad ISO ↔ Requisitos}
\begin{longtable}{p{3cm}p{3cm}p{8cm}}
\toprule
\textbf{Norma ISO} & \textbf{Capítulo} & \textbf{Requisito aplicado} \\
\midrule
ISO 15489 & 7–9 & Procesos del ciclo de vida del documento \\
ISO 30301 & 4–10 & Sistema de gestión y auditoría \\
ISO 23081 & 5–7 & Estructura y gestión de metadatos \\
ISO 16175 & 6–9 & Funcionalidades del sistema (captura, clasificación, auditoría) \\
ISO 26102 & 5–8 & Requisitos técnicos, trazabilidad y conformidad \\
\bottomrule
\end{longtable}

\end{document}
